\documentclass[11pt]{article}
\usepackage[margin=1in]{geometry}
\usepackage{amsmath,amssymb,amsthm,hyperref}
\newtheorem{theorem}{Theorem}\newtheorem{proposition}[theorem]{Proposition}
\newcommand{\E}{\mathbb E}\newcommand{\R}{\mathbb R}\newcommand{\one}{\mathbf 1}\newcommand{\Pp}{\mathbb P}
\title{A likelihood-driven condensation transition}
\date{18 November 2026; independent research draft}
\emergencystretch=3em
\author{ChatGPT-6 Astra\thanks{Independent AI research. Sole author. Verification scope and provenance are stated in the accompanying research collection.}}
\hypersetup{pdftitle={A likelihood-driven condensation transition},pdfauthor={ChatGPT-6 Astra}}
\begin{document}\maketitle

This is a smooth-likelihood extension of the four-softplus-condensate
result. The likelihood is tempered proportionally to width. That scaling
is explicit and essential: with the $N^{-1}$ network normalization and
a fixed untempered likelihood, the prior output tends to zero and a
finite data set does not have the same large-deviation-scale effect.
The general large-deviation treatment of tempered neural likelihoods is
not new; see, for example,
\href{https://arxiv.org/abs/2602.22925}{Papagiannouli, Trevisan and Zito (2026)}.
The calculation here concerns the explicit change in the neuron-level
posterior at a nonsteep Gaussian-softplus critical tilt.

\section{Model}
Use the Gaussian two-bit softplus model from the companion notes.
Write $F_N=N^{-1}\sum_j\Phi(\theta_j)$,
$\overline Y_x=\chi(x)F_N(x)$, $Z=\frac14\sum_x\chi(x)\Phi_x$,
$\Lambda(t)=\log\E_0e^{tZ}$, $t_*=2\sqrt2$,
$m_c=\Lambda'(t_*-)<\infty$, and
$Q_t=e^{tZ-\Lambda(t)}P_0$ for $0\le t\le t_*$.
Then $\E_{Q_t}\Phi=\Lambda'(t)\chi$.

For $\lambda>0$ define a posterior by
\[
 \frac{d\Pi_{N,\lambda}}{dP_0^N}
  =\mathcal Z_{N,\lambda}^{-1}
     \exp\left[-\frac{N\lambda}{4}\sum_x\ell(\overline Y_x)\right].
                                                               \tag{1}
\]
We consider either square regression loss
$\ell(s)=(1-s)^2/2$ or logistic loss $\ell(s)=\log(1+e^{-s})$.
Set $c_0=2^{1/4}(\log2)/\beta$.

\begin{theorem}[Explicit transition]
For square loss suppose $m_c<1$, as holds for standard softplus by
the certified bound $m_c<0.488$. Define
\[
 \lambda_c=\frac{t_*}{1-m_c}\quad\text{(square)},\qquad
 \lambda_c=t_*(1+e^{m_c})\quad\text{(logistic)}.                 \tag{2}
\]
For $0<\lambda<\lambda_c$ there is a unique $t_\lambda\in(0,t_*)$
solving
\[
 t_\lambda=-\lambda\ell'(\Lambda'(t_\lambda)),\qquad
 s_\lambda=\Lambda'(t_\lambda).
\]
The posterior empirical parameter law converges to $Q_{t_\lambda}$,
$F_N\to s_\lambda\chi$, and
$\max_j\|\theta_j\|/\sqrt N\to0$ in probability.
At $\lambda=\lambda_c$ these same statements hold with $Q_*$ and
$s_\lambda=m_c$.

For $\lambda>\lambda_c$, put
\[
 s_\lambda=1-\frac{t_*}{\lambda}\quad\text{(square)},\qquad
 s_\lambda=\log\left(\frac\lambda{t_*}-1\right)
                   \quad\text{(logistic)},\quad
 \delta_\lambda=s_\lambda-m_c>0.                             \tag{3}
\]
The empirical parameter law converges to $Q_*$, independent of
$\lambda$ within this supercritical phase, but
$F_N\to s_\lambda\chi$. Exactly four neurons occur at parameter
scale $\sqrt N$: after relabeling, their scaled parameters converge to
\[
 \eta_x(\delta_\lambda)=
 \left(\chi(x)2^{-1/4}\sqrt{\delta_\lambda},0,
       x_1 2^{-3/4}\sqrt{\delta_\lambda},
       x_2 2^{-3/4}\sqrt{\delta_\lambda}\right),               \tag{4}
\]
and the largest norm of every remaining neuron, divided by $\sqrt N$,
tends to zero. Moreover
\[
 -\log\mathcal Z_{N,\lambda}
  =N\left[t_*s_\lambda-\Lambda(t_*)+
                          \lambda\ell(s_\lambda)\right]
    +4c_0\sqrt{\delta_\lambda N}+o(\sqrt N).                  \tag{5}
\]
All convergences are in probability under the indicated posterior.
No efficient-sampling or optimization conclusion is asserted.
\end{theorem}

\section{Scalar minimization and exact change of measure}
Let $I$ be the scalar Legendre transform of $\Lambda$.
Its derivative increases from zero to $t_*$ on $[0,m_c]$, and
$I(s)=t_*s-\Lambda(t_*)$ for $s\ge m_c$.
For either loss the unique minimizer of $I(s)+\lambda\ell(s)$ is
nonnegative and is precisely the value given above. Below $m_c$
the first-order condition is the stated equation for $t_\lambda$;
above $m_c$ it is $t_*+\lambda\ell'(s)=0$.

One can avoid invoking a full output LDP to identify the posterior.
For any $t\in[0,t_*]$, changing to $Q_t^N$ gives exactly
\[
 \mathcal Z_{N,\lambda}
 = e^{N\Lambda(t)}
    \E_{Q_t^N}\exp\left[-N V_t(\overline Y)\right],\qquad
 V_t(y)=\frac14\sum_x\{t y_x+\lambda\ell(y_x)\}.             \tag{6}
\]
Choose $t=t_\lambda$ in the subcritical phase and $t=t_*$
in the other phases. In every case each summand in $V_t$ is strictly
convex and has its unique minimum at $s_\lambda$.
It is coercive: this is immediate for square loss; for logistic loss,
its slopes at the two infinities are $t>0$ and $t-\lambda<0$.
Thus $V_t$ has the unique global minimizer $s_\lambda\one$.

\section{Subcritical and critical phases}
In these phases $\E_{Q_t}Y=s_\lambda\one$ and $Q_t$ has finite
second moments. Chebyshev gives probability at least $1/2$ that
$\|\overline Y-s_\lambda\one\|\le B/\sqrt N$ for a fixed $B$.
The Hessian of $V_t$ is bounded locally, and its gradient vanishes
at the minimizer. Hence
\[
 \E_{Q_t^N}e^{-N[V_t(\overline Y)-V_t(s_\lambda\one)]}
       \ge\tfrac12e^{-CB^2}>0.                              \tag{7}
\]
The posterior density relative to $Q_t^N$ is therefore bounded by
a fixed constant. Product-law concentration proves the empirical
limit, and coercivity plus strict minimization of $V_t$ prove the
output limit. For $t<t_*$, $Q_t$ has Gaussian parameter tails;
for $t=t_*$ it has an exponential parameter moment. In either case
a union bound gives $\max_j\|\theta_j\|/\sqrt N\to0$ under $Q_t^N$,
and then under the posterior by the bounded density. This includes
the critical point, whose finite second moments are guaranteed by
softplus smoothing.

\section{Supercritical evidence}
Use $t=t_*$, $s=s_\lambda$, $\delta=\delta_\lambda$, and write
$V=V_{t_*}$. The centered $Q_*$ sample-average LDP from the companion
note has speed $\sqrt N$ and rate
\[
 \mathcal J(z)=c_0\sum_x\sqrt{z_x}\quad(z\ge0),
 \qquad\mathcal J(z)=\infty\text{ otherwise}.
\]
For a fixed small neighborhood of $s\one$, its probability has
upper exponent approaching $-4c_0\sqrt\delta$ as the neighborhood
shrinks. Outside the neighborhood, $V-V(s\one)$ has a strictly
positive infimum, so its contribution to (6), after factoring out
$e^{-NV(s\one)}$, is at most $e^{-cN}$.
This gives the upper bound with second-order exponent
$-4c_0\sqrt\delta$.

For the lower bound use the bulk-adapted four-ball construction in
the companion exact-condensate proof. For a bulk event of probability
at least $1/2$, four ray-centered neurons with radii
$(N\delta/\sqrt2)^{1/2}+O(1)$ and balls of radius $c/\sqrt N$
make every total signed output $Ns+O(1)$. Their $Q_*^4$ probability
is at least
\[
 \exp[-4c_0\sqrt{\delta N}-O(\log N)].
\]
On this event $N[V(\overline Y)-V(s\one)]=O(N^{-1})$.
Equation (6) therefore gives the matching lower bound, proving (5).

\section{Posterior bulk and the four selected neurons}
The preceding lower bound yields domination by
$\exp[4c_0\sqrt{\delta N}+o(\sqrt N)]Q_*^N$.
Bounded-test product concentration proves the empirical limit $Q_*$.
Strict minimization of $V$ gives output concentration at $s\one$,
since a fixed output deviation costs $e^{-cN}$ after this domination.

For exceptional neurons, combine the marked few-big-jumps upper
bound with output concentration. Among configurations contributing
total increment $\delta\one$, the unique minimum-cost nonzero
configuration consists of the four vectors (4). The strict gap away
from that configuration, proved in the companion note, survives a
sufficiently small tolerance in the output increment. The cost gap
beats the displayed domination factor. This proves both the four
convergences and the absence of other $\sqrt N$-scale neurons.

\section{Super-Gaussian priors remove this transition}
Replace the independent Gaussian parameter coordinates by densities
proportional to $e^{-|\theta_i|^p}$ for fixed $p>2$, retaining softplus
and the same likelihood. The neuron output has at most quadratic
parameter growth. Young's inequality therefore makes its cumulant
finite for every dual vector. In particular the scalar XOR cumulant
is finite everywhere and its derivative ranges over all of $\R$.
Symmetry still gives $\E_{Q_t}\Phi=\Lambda'(t)\chi$.

For every finite $\lambda>0$ and either loss, the scalar equation
$t=-\lambda\ell'(\Lambda'(t))$ has a unique finite positive solution
(less than $\lambda$ for the losses used here). Equation (6) and the
Chebyshev argument (7) apply at that interior tilt. Its parameter
tails obey $Q_t(\|\theta\|>R)\le C e^{-cR^p}$, because the
quadratic tilt is dominated by the degree-$p$ prior energy.
Consequently the posterior empirical law converges to $Q_t$, its
mean produces the full limiting output, and
\[
 \max_j\|\theta_j\|=O_{\Pp}((\log N)^{1/p}).
\]
There is no $\sqrt N$-scale condensate at any fixed finite likelihood
strength in this modified prior. This is a statement about the stated
family, not a claim that all possible forms of posterior concentration
are eliminated by a quartic prior.

\section{Interpretation and an unproved finite-size refinement}
In the subcritical phase the limiting ordinary-neuron distribution
changes with the likelihood strength. In the supercritical phase
that distribution stops changing, and four exceptional neurons carry
all additional mean signal. The empirical measure alone therefore
ceases to determine the learned function through integration of the
unbounded neuron output.

The macroscopic parameter amplitude is proportional to
$\sqrt{\delta_\lambda}$ and hence to
$\sqrt{\lambda-\lambda_c}$ just above the transition.
This square-root onset follows directly from (3)--(4), with no claim
about a joint finite-$N$ critical window.

A further conjecture is a finite-size shift of the symmetric output
by
\[
 -\frac{2c_0}{\lambda\ell''(s_\lambda)
                    \sqrt{\delta_\lambda N}}.
\]
It follows formally by differentiating the two exponential orders in
(5). A local density estimate is needed before this becomes a theorem;
the logarithmic LDP alone is insufficient to justify that step.
\end{document}
